% !TEX root = ../../main.tex

\renewcommand{\thesection}{2.1}

\section{行列式的定义}\label{sec:2.1}

\begin{ex}{3}\label{ex:2.1.3}
    计算行列式
    \[
    \begin{vmatrix}
        0 & 1 & 0 & \cdots & 0\\
        0 & 0 & 2 & \cdots & 0\\
        \vdots & \vdots & \vdots & \ddots & \vdots\\
        0 & 0 & 0 & \cdots & n-1\\
        n & 0 & 0 & \cdots & 0
    \end{vmatrix}.
    \]
\end{ex}

\begin{solution}
    设待求 \(n\) 阶行列式为 \(\det(a_{ij})\). 由行列式定义, 有
    \[
    \det(a_{ij})=(-1)^{\tau(234\cdots n1)}a_{12}a_{23}a_{34}\cdots a_{(n-1)n}a_{n1}
    =(-1)^{n-1}n!.\qedhere
    \]
\end{solution}

\begin{ex}{5}\label{ex:2.1.5}
    计算
    \[
    f(x)=
    \begin{vmatrix}
        x & x & 2 & -1\\
        2 & 3x & 0 & 2\\
        3 & 2 & x & 1\\
        1 & 1 & 1 & x\\
    \end{vmatrix}
    \]
    中 \(x^4\) 和 \(x^3\) 的系数.
\end{ex}

\begin{solution}
    设题中行列式为 \(\det(a_{ij})\). 在其展开式中的所有项中, 次数为 \(4\) 的项之和为
    \[
    (-1)^{\tau(1234)}a_{11}a_{22}a_{33}a_{44}=3x^4,
    \]
    次数为 \(3\) 的项之和为
    \[
    (-1)^{\tau(2134)}a_{12}a_{21}a_{33}a_{44}=-2x^3.
    \]
    
    因此, \(x^4\) 的系数为 \(3\), \(x^3\) 的系数为 \(-2\).
\end{solution}

\begin{ex}{6}\label{ex:2.1.6}
    证明
    \[
    \begin{vmatrix}
        0 & & a_1\\
        & \iddots &\\
        a_n & & *
    \end{vmatrix}
    =
    \begin{vmatrix}
        * & & a_1\\
        & \iddots &\\
        a_n & & 0
    \end{vmatrix}
    =(-1)^{n(n-1)/2}a_1a_2\cdots a_n.
    \]
\end{ex}

\begin{proof}
    设题中行列式分别为 \(\det(b_{ij}),\det(c_{ij})\). 由行列式定义, 有
    \begin{gather*}
        \det(b_{ij})=\sum_{p_1p_2\cdots p_n\in S_n}(-1)^{\tau(p_1p_2\cdots p_n)}\prod_{i=1}^{n}b_{ip_i},\\
        \det(c_{ij})=\sum_{p_1p_2\cdots p_n\in S_n}(-1)^{\tau(p_1p_2\cdots p_n)}\prod_{i=1}^{n}c_{ip_i},
    \end{gather*}
    其中 \(S_n\) 是所有 \(n\) 排列之集.

    考虑其中所有不为 \(0\) 的项, 有
    \begin{align*}
        \det(b_{ij})=\det(c_{ij})&=(-1)^{\tau(n(n-1)\cdots1)}\prod_{i=1}^{n}a_i\\
        &=(-1)^{n-1+n-2+\cdots+0}\prod_{i=1}^{n}a_i\\
        &=(-1)^{n(n-1)/2}\prod_{i=1}^{n}a_i.\qedhere
    \end{align*}
\end{proof}